% ============================================================
% 3.2.3.3 Inflation and Deflation — Mark Scheme
% AQA AS Economics
% ============================================================

\documentclass[12pt,a4paper]{article}
\usepackage[margin=1in,headheight=15pt]{geometry}
\usepackage{enumitem}
\usepackage{amsmath}
\usepackage{fancyhdr}
\usepackage{booktabs}
\usepackage{array}
\RequirePackage{tcolorbox}
\tcbuselibrary{skins}

% Key Point — yellow; for the single most important idea on a slide
\newtcolorbox{keypoint}{
    enhanced,
    colback=yellow!12,
    colframe=yellow!60!black,
    fonttitle=\bfseries,
    title=Key Point,
    arc=2pt,
    boxrule=0.8pt,
}

% Formula — blue; for named equations and their variable key
\newtcolorbox{formula}[1][Formula]{
    enhanced,
    colback=blue!5,
    colframe=blue!45!black,
    fonttitle=\bfseries,
    title=#1,
    arc=2pt,
    boxrule=0.8pt,
}

% Example Question — teal; for exam-style questions posed to students
\newtcolorbox{examplequestion}[1][Example Question]{
    enhanced,
    colback=teal!6,
    colframe=teal!55!black,
    fonttitle=\bfseries,
    title=#1,
    arc=2pt,
    boxrule=0.8pt,
}

% Example Solution — green; always follows an examplequestion block
\newtcolorbox{examplesolution}[1][Solution]{
    enhanced,
    colback=green!6,
    colframe=green!45!black,
    fonttitle=\bfseries,
    title=#1,
    arc=2pt,
    boxrule=0.8pt,
}

% ---- Worksheet environments --------------------------------

% Scenario — blue; for group discussion prompts on worksheets
% Usage: \begin{scenario}{Scenario 1: Title with commas, fine} ... \end{scenario}
\newtcolorbox{scenario}[1]{
    enhanced,
    colback=blue!5,
    colframe=blue!45!black,
    fonttitle=\bfseries,
    title={#1},
    arc=2pt,
    boxrule=0.8pt,
}

% Instructions — blue; for worksheet-level instruction blocks
\newtcolorbox{instructions}{
    enhanced,
    colback=blue!5,
    colframe=blue!45!black,
    fonttitle=\bfseries,
    title=Instructions,
    arc=2pt,
    boxrule=0.8pt,
}

% Extract — grey; for data response source material
% Usage: \begin{extract}{Source: Author, Year} ... \end{extract}
\newtcolorbox{extract}[1]{
    enhanced,
    colback=gray!8,
    colframe=gray!50,
    fonttitle=\bfseries,
    title={#1},
    arc=2pt,
    boxrule=0.8pt,
}


\pagestyle{fancy}
\fancyhf{}
\lhead{Dr Jac Thomas \quad\textit{Mark Scheme}}
\rhead{Topic index: 3.2.3.3}
\rfoot{Page \thepage}

% Award marks generously for equivalent correct answers.
% Do not penalise candidates for poor spelling or grammar
% unless it obscures the economic meaning.

\begin{document}

\begin{center}
    \Large\textbf{Inflation and Deflation --- Mark Scheme}\\[5pt]
    \normalsize Topic index: 3.2.3.3 \quad|\quad AQA AS Economics\\[3pt]
    \textit{Total Marks: 60}
\end{center}

\vspace{1em}

% ============================================================
\section*{Section A: Key Terms (10 marks)}
% ============================================================

\textit{Award 1 mark for a basic correct statement; 2 marks for a precise, complete definition.}

\vspace{0.5em}

\begin{enumerate}[leftmargin=*, label=\textbf{\arabic*.}]

    \item \textbf{Inflation} \hfill (2 marks)\\
    A sustained rise in the general price level of goods and services in an economy over time.\\
    \textit{1 mark: prices rising. 2 marks: must include ``sustained'' / ``general'' (not just one good) or equivalent.}

    \vspace{0.5em}

    \item \textbf{Deflation} \hfill (2 marks)\\
    A sustained fall in the general price level, i.e.\ a negative rate of inflation.\\
    \textit{1 mark: prices falling. 2 marks: must specify ``sustained'' / ``general price level'' falling, or state the inflation rate is negative. Do not award 2 marks if candidate confuses deflation with disinflation.}

    \vspace{0.5em}

    \item \textbf{Disinflation} \hfill (2 marks)\\
    A fall in the rate of inflation; the price level is still rising but at a slower rate than before.\\
    \textit{1 mark: rate of inflation falling. 2 marks: must make clear that prices are still rising (not falling) --- the rate of increase is slowing, not reversing. Award 0 marks if candidate defines disinflation as falling prices (that is deflation).}

    \vspace{0.5em}

    \item \textbf{Demand-pull inflation} \hfill (2 marks)\\
    Inflation caused by a rise in aggregate demand, particularly when the economy is operating at or near full capacity, pulling the price level upward.\\
    \textit{1 mark: links inflation to rising AD. 2 marks: must reference full capacity / positive output gap, or explain that AD rises faster than productive capacity.}

    \vspace{0.5em}

    \item \textbf{Cost-push inflation} \hfill (2 marks)\\
    Inflation caused by a rise in the costs of production (e.g.\ wages, raw materials, energy), which shifts the short-run aggregate supply curve leftward, raising the price level while reducing output.\\
    \textit{1 mark: links inflation to rising costs / SRAS shift. 2 marks: must identify a specific cost factor or explain the SRAS mechanism. Accept any valid cost factor.}

\end{enumerate}

\newpage

% ============================================================
\section*{Section B: Group Discussion --- Indicative Answers}
% ============================================================

\textit{No marks awarded here. Use these notes to guide class feedback.}

\vspace{0.5em}

\textbf{Scenario 1: UK Energy Crisis (2021--22)}
\begin{itemize}
    \item Type: \textbf{Cost-push inflation} --- rising global gas prices increased production costs across manufacturing, transport, and food industries.
    \item Cause: Supply-side shock --- global gas supply disruptions (following the Russia--Ukraine war and post-COVID demand recovery) drove up energy prices.
    \item Policy: Difficult to address with conventional monetary policy (raising interest rates does not lower gas prices). Government used targeted support (e.g.\ Energy Price Guarantee, cost-of-living payments) to cushion the blow. The Bank of England raised rates to prevent the cost-push shock from becoming embedded in wage and price expectations.
\end{itemize}

\vspace{0.5em}

\textbf{Scenario 2: Post-COVID Demand Surge (2021)}
\begin{itemize}
    \item Type: \textbf{Demand-pull inflation} (with a supply-chain element of cost-push). Primarily driven by a surge in consumer spending.
    \item Cause: Accumulated household savings were released as restrictions lifted; pent-up demand exceeded supply capacity, pushing prices up.
    \item Policy: Contractionary monetary policy (higher interest rates) to reduce borrowing and cool consumer spending; reduce the rate of AD growth back toward the economy's productive capacity.
\end{itemize}

\vspace{0.5em}

\textbf{Scenario 3: Japan's Lost Decade}
\begin{itemize}
    \item Illustrates: \textbf{Deflation} and the risk of a deflationary spiral.
    \item Why falling prices are harmful: Consumers and firms delay spending, expecting prices to fall further; this reduces AD, causing further falls in output and employment; the real burden of debt rises; investment collapses. A self-reinforcing cycle of decline.
    \item Deflation vs disinflation: Deflation = price level falling (negative inflation). Disinflation = price level still rising but more slowly. Japan experienced actual deflation --- the price level fell.
    \item Policy to escape: Expansionary fiscal policy (large government spending increases); unconventional monetary policy (quantitative easing, negative interest rates); structural reform to boost productivity and long-run aggregate supply.
\end{itemize}

\newpage

% ============================================================
\section*{Section C: Multiple Choice (10 marks, 1 each)}
% ============================================================

\vspace{0.5em}

\begin{table}[h]
    \centering
    \renewcommand{\arraystretch}{1.3}
    \begin{tabular}{@{}ccc@{}}
        \toprule
        \textbf{Question} & \textbf{Answer} & \textbf{Correct option} \\ \midrule
        1  & B & The rate of inflation is falling while prices are still rising \\
        2  & B & Consumer spending surges as the economy approaches full employment \\
        3  & C & A leftward shift of the SRAS curve \\
        4  & B & Changes in the weighted average price of a basket of goods and services \\
        5  & A & Raising production costs, shifting SRAS leftward \\
        6  & C & A rise in firms' wage costs following a large increase in the National Living Wage \\
        7  & B & Disinflation \\
        8  & C & Reduce the growth of aggregate demand in order to reduce inflation \\
        9  & B & Increasing demand for UK exports, raising AD and putting upward pressure on prices \\
        10 & B & Cost-push inflation from rising input costs shifting SRAS leftward \\ \bottomrule
    \end{tabular}
\end{table}

\newpage

% ============================================================
\section*{Section D: Quantitative Skills (8 marks)}
% ============================================================

\begin{enumerate}[leftmargin=*, label=\textbf{\arabic*.}]

    \item \textbf{Food price index for 2023} \hfill (2 marks)
    \[
        \text{Price index} = \frac{\text{Price in current year}}{\text{Price in base year}} \times 100
        = \frac{4.80}{4.00} \times 100 = \mathbf{120}
    \]
    \textit{1 mark: correct formula / method. 1 mark: correct answer (120).}

    \vspace{0.5em}

    \item \textbf{Transport price index for 2023} \hfill (1 mark)
    \[
        \frac{64.80}{60.00} \times 100 = \mathbf{108}
    \]
    \textit{1 mark: 108.}

    \vspace{0.5em}

    \item \textbf{Overall CPI for Econland in 2023} \hfill (2 marks)
    \[
        \text{CPI} = \frac{25 \times 120 + 30 \times 110 + 25 \times 108 + 20 \times 105}{100}
        = \frac{3000 + 3300 + 2700 + 2100}{100}
        = \frac{11{,}100}{100} = \mathbf{111.0}
    \]
    \textit{1 mark: correct method (weighted average using all four indices). 1 mark: correct answer (111.0). Award follow-through marks if Q1 or Q2 answers are incorrect but method is correct.}

    \vspace{0.5em}

    \item \textbf{Rate of inflation, 2022--2023} \hfill (1 mark)
    \[
        \text{Inflation rate} = \frac{\text{CPI}_t - \text{CPI}_{t-1}}{\text{CPI}_{t-1}} \times 100
        = \frac{111.0 - 100}{100} \times 100 = \mathbf{11.0\%}
    \]
    \textit{1 mark: 11.0\% (or 11\%). Accept follow-through from Q3.}

    \vspace{0.5em}

    \item \textbf{Inflation, deflation, or disinflation?} \hfill (2 marks)\\
    This represents \textbf{disinflation}. The rate of inflation has fallen (from 11\% to 3\%), but prices are still rising --- the price level has not fallen. Deflation would require a \textit{negative} inflation rate (i.e.\ the price level actually falling). A 3\% inflation rate means prices are still increasing, just more slowly than before.\\
    \textit{1 mark: correctly identifies disinflation. 1 mark: explains that prices are still rising (not falling) --- the rate of increase has slowed.}

\end{enumerate}

\newpage

% ============================================================
\section*{Section E: Diagrams (10 marks)}
% ============================================================

\textit{Mark for: correct shape of curves, correct labelling, correct equilibrium position(s), and correct identification of the shift and its effect on the price level.}

\vspace{0.5em}

\textbf{Question 1 --- Demand-pull inflation} \hfill (5 marks)

\begin{itemize}
    \item \textbf{[1]} Downward-sloping AD curve, upward-sloping SRAS curve --- both labelled. (LRAS optional but accept if present.)
    \item \textbf{[1]} Second AD curve (AD$_2$) drawn to the \textit{right} of AD$_1$ --- shift direction correct. Arrow indicating rightward shift or AD$_2$ clearly labelled.
    \item \textbf{[1]} New equilibrium $E_2$ at the intersection of AD$_2$ and SRAS, labelled.
    \item \textbf{[1]} New price level $P_2 > P_1$ correctly labelled on the y-axis.
    \item \textbf{[1]} New real GDP $Y_2 > Y_1$ correctly labelled on the x-axis. If LRAS is drawn, award this mark for showing $Y_2$ is at or beyond $Y_f$ (positive output gap).
\end{itemize}

\vspace{0.8em}

\textbf{Question 2 --- Cost-push inflation} \hfill (5 marks)

\begin{itemize}
    \item \textbf{[1]} Downward-sloping AD curve and upward-sloping SRAS curve --- both labelled. Initial equilibrium $E_1$ shown.
    \item \textbf{[1]} Second SRAS curve (SRAS$_2$) drawn to the \textit{left} (and above) of SRAS$_1$ --- shift direction correct. Arrow indicating leftward shift or SRAS$_2$ clearly labelled.
    \item \textbf{[1]} New equilibrium $E_2$ at the intersection of SRAS$_2$ and AD, labelled.
    \item \textbf{[1]} New price level $P_2 > P_1$ correctly labelled on the y-axis.
    \item \textbf{[1]} New real GDP $Y_2 < Y_1$ correctly labelled on the x-axis (output \textit{falls} --- this distinguishes cost-push from demand-pull). Penalise if student shows output rising in Q2.
\end{itemize}

\newpage

% ============================================================
\section*{Section F: Short Response (7 marks)}
% ============================================================

\begin{enumerate}[leftmargin=*, label=\textbf{\arabic*.}]

    \item \textbf{Difference between deflation and disinflation} \hfill (4 marks)\\
    \textit{Award 2 marks per point: 1 mark for stating the difference, 1 mark for developing it.}

    \vspace{0.3em}
    Indicative answer:
    \begin{itemize}
        \item \textbf{Deflation:} The price level is \textit{falling} --- the rate of inflation is negative. In deflation, the general level of prices in the economy is lower than it was in the previous period. This can be harmful because it encourages consumers to delay spending (waiting for further price falls), potentially causing a deflationary spiral.
        \item \textbf{Disinflation:} The rate of inflation is \textit{falling} but remains positive --- prices are still rising, just more slowly. For example, if inflation falls from 8\% to 3\%, prices are still higher than the previous year; they are simply rising at a slower rate. Disinflation is not the same as a falling price level.
    \end{itemize}
    \textit{Award up to 4 marks for a clear explanation of both terms including what happens to the price level in each case. A response that confuses the two (e.g.\ states disinflation means prices fall) cannot score more than 2 marks.}

    \vspace{0.5em}

    \item \textbf{One way a rise in world commodity prices could lead to higher UK inflation} \hfill (3 marks)\\
    \textit{1 mark: identify a relevant commodity or transmission mechanism. 2 marks: explain the mechanism clearly.}

    \vspace{0.3em}
    Indicative answers (any one):
    \begin{itemize}
        \item \textbf{Rise in oil prices:} Oil is used as an input across many sectors --- fuel, plastics, chemicals, transport. A rise in world oil prices increases firms' production costs. Firms pass these higher costs on to consumers in the form of higher prices. In AD/AS terms, SRAS shifts leftward, raising the price level --- this is cost-push inflation.
        \item \textbf{Rise in food commodity prices (wheat, soya):} The UK imports a significant share of its food. If global wheat prices rise (e.g.\ due to drought or conflict disrupting supply), UK food manufacturers and retailers face higher input costs. These are passed on in higher food prices, raising the CPI.
        \item \textbf{Rise in metals/raw material prices:} Manufacturing firms face higher input costs; factory-gate prices rise; these are passed on throughout the supply chain, raising the general price level.
    \end{itemize}

\end{enumerate}

\newpage

% ============================================================
\section*{Section G: Essay (15 marks)}
% ============================================================

\textit{``Evaluate the view that rising aggregate demand is the main cause of inflation in the UK economy.''}

\vspace{0.5em}

\begin{extract}{Marking guidance}
    Level 4 (13--15): Thorough, balanced analysis of both demand-pull and cost-push causes. Strong evaluation --- considers when demand-pull is most likely (economy near $Y_f$, positive output gap) versus when cost-push dominates (supply shocks, imported inflation). Uses real-world evidence effectively. Well-structured and precise.\\[0.4em]
    Level 3 (9--12): Good analysis of demand-pull and at least one alternative cause. Some evaluation of which is ``main'' cause and under what conditions. May lack depth in evaluation or real-world application.\\[0.4em]
    Level 2 (5--8): Some relevant analysis, but limited engagement with cost-push or global factors. Evaluation may be asserted rather than developed. May focus too heavily on one cause.\\[0.4em]
    Level 1 (1--4): Basic or largely descriptive answer. May only define inflation or list causes without analysis. Little or no evaluation.
\end{extract}

\vspace{0.8em}

\textbf{Indicative content}

\vspace{0.3em}
\textbf{Arguments that demand-pull is the main cause:}
\begin{itemize}
    \item Rising AD (from cuts in interest rates, higher consumer confidence, expansionary fiscal policy) creates inflationary pressure, particularly when the economy is near full employment ($Y_e \approx Y_f$).
    \item The post-COVID episode (2021): surge in consumer spending as savings were released; pent-up demand outstripped supply capacity, contributing to rising prices.
    \item Strong global growth boosts UK export demand, raising AD and adding demand-side inflationary pressure.
    \item The Bank of England's remit focuses on managing demand --- the existence of the 2\% target and interest-rate tool reflects the importance of the demand side.
\end{itemize}

\vspace{0.3em}
\textbf{Arguments against (cost-push and global causes matter too):}
\begin{itemize}
    \item \textbf{Cost-push inflation} arises independently of domestic demand: a rise in oil prices, global food shortages, or a depreciation of sterling can push SRAS leftward even when the economy is operating below full capacity.
    \item The 2021--22 UK inflation surge (reaching 11.1\%) was primarily supply-driven: global energy prices, supply chain disruption, and the Russia--Ukraine war raised input costs across the economy --- AD was not the dominant driver.
    \item \textbf{Imported inflation}: the UK is a price-taker in global commodity markets. Even if UK domestic demand is stable, higher world oil or food prices can raise UK inflation via cost-push channels.
    \item \textbf{Exchange rate}: a depreciation of sterling raises import costs, feeding into SRAS regardless of the level of domestic AD.
\end{itemize}

\vspace{0.3em}
\textbf{Evaluation:}
\begin{itemize}
    \item The dominant cause depends on the state of the economy and the prevailing shocks: demand-pull inflation is most likely when the economy is near $Y_f$ (positive or closing output gap); cost-push can occur even in recessions (stagflation).
    \item In the UK context, the 2021--23 inflation was widely attributed mainly to supply shocks (energy, supply chains), not excess demand --- suggesting demand-pull is \textit{not always} the main cause.
    \item In practice both causes often occur simultaneously and reinforce each other (e.g.\ supply-side cost increases raise inflation expectations, leading to wage demands that shift SRAS further).
    \item \textbf{Conclusion:} Rising AD is an important cause of inflation, particularly in a booming economy near full employment. However, cost-push factors --- especially imported commodity prices and exchange rate movements --- can be equally or more important, as the 2022 episode illustrates. The statement is therefore partially correct but needs qualification.
\end{itemize}

\vfill
\begin{center}
    \textbf{[Total: 60 marks]}
\end{center}

\end{document}
